\documentclass[12pt]{article}
\usepackage[left=1.2in,right=1.2in,top=1.2in,bottom=1.2in]{geometry}
\usepackage{lmodern}
\usepackage[T1]{fontenc}
\usepackage{microtype}
\usepackage{setspace}
\usepackage{parskip}
\usepackage[hidelinks]{hyperref}

\pagestyle{empty}
\setlength{\parindent}{0pt}
\setlength{\parskip}{1em}

\begin{document}

\singlespacing

Tingo Mar\'ia, April 13, 2026

\medskip

\textbf{Editorial Board}\\
\textit{Resources and Conservation Recycling}\\
Elsevier

\medskip
\noindent\rule{\linewidth}{0.4pt}
\medskip

Dear Editor,

We submit for your consideration the manuscript titled \textbf{``Municipal Public Cleaning Expenditure and Solid Waste Management: Evidence from Peru, 2015--2019''} by Dante Barreto Gamarra (Universidad Nacional Agraria de la Selva, UNAS).

\textbf{Research question and motivation.} Despite a growing consensus that fiscal capacity constrains solid waste management in low- and middle-income countries, rigorous causal evidence at the sub-national level remains scarce. The core challenge is endogeneity: municipalities with stronger institutional capacity simultaneously spend more and report better outcomes, making standard panel estimates uninformative for policy. This paper directly addresses that gap using a nationally representative municipal panel from Peru.

\textbf{Contribution.} We estimate the causal effect of municipal public cleaning expenditure on three dimensions of solid waste management --- daily waste collection volume, a planning adoption index (PIGARS), and the probability of formal disposal in sanitary landfills --- across 1,874 Peruvian municipalities from 2015 to 2019 (9,322 observations; 196 provincial clusters). To correct for endogeneity, we apply the Correlated Random Effects with Control Function estimator (CRE+CF) proposed by Wooldridge (2015), combined with the Mundlak (1978) device. This design controls for time-invariant unobservables while purging time-varying endogeneity through first-stage residuals, addressing the main identification threat in municipal public finance studies.

\textbf{Key findings.} A 10\% increase in accrued expenditure raises daily waste collection by 8.79\% (elasticity $= +0.879$; $p < 0.001$), increases the planning index by 0.030 points on a five-instrument scale ($+0.300$ per log-unit; $p < 0.001$), and raises the probability of formal sanitary landfill disposal by 2.2 percentage points over an 18\% baseline (AME $= +0.022$; $p < 0.001$). All three primary hypotheses survive Bonferroni-Holm correction. Robustness checks include Oster (2019) omitted variable bounds ($\delta^* < 0$ for all outcomes, $|\delta^*| \geq 1.35$), DL(1) pre-trend tests (lag ratios $< 0.10$), and Conley--Hansen--Rossi (2012) local-to-exogenous sensitivity analysis.

\textbf{Fit for this journal.} The paper speaks directly to the empirical literature on municipal solid waste management costs and governance in developing countries --- a line of research that this journal has helped define (see, e.g., Bel \& Fageda, 2010, \textit{RCR} 54; Guerrero et al., 2013, \textit{Waste Management} 33). We extend that literature by providing identified causal estimates --- rather than cost-function associations --- and by documenting the mechanism through which budgetary constraints bind the transition to formal disposal infrastructure, a policy-relevant finding for Peru's Decreto Legislativo N\textdegree~1278 and analogous LMIC frameworks.

\textbf{Originality statement.} This manuscript has not been published elsewhere and is not under simultaneous review at any other journal.

We believe this paper makes a methodological and substantive contribution well suited to your readership. We look forward to the editorial assessment.

\medskip

Sincerely,

\bigskip\bigskip

\textbf{Dante Barreto Gamarra}\\
Faculty of Economics\\
Universidad Nacional Agraria de la Selva (UNAS)\\
Tingo Mar\'ia, Per\'u

\end{document}
